% Section: Introduction
\section{Introduction}
\label{sec:introduction}

% TODO: Expand this section.
%
% Key topics:
% - Motivation for reflective development systems
% - The challenge of AI-augmented recursive engineering
% - Overview of the Reflector framework
% - Contributions of this paper
% - Paper organization

The emergence of large language model (LLM)-powered development agents has fundamentally
altered the relationship between human engineers and software systems. Where traditional
development workflows position humans as the primary actors—writing, reviewing, and
deploying code—modern AI-augmented pipelines increasingly delegate these activities to
autonomous agents capable of operating across extended time horizons with minimal
per-step human oversight.

This shift introduces substantial productivity benefits: agents can generate, test, and
refine code faster than human developers, maintain consistency across large codebases,
and parallelize work across multiple concurrent tasks. However, it also introduces a
new class of systemic risks that existing software engineering methodologies are
ill-equipped to address.

Chief among these risks is \textit{recursive drift}: the progressive divergence of an
autonomous system from its intended objectives as it operates through self-referential
feedback cycles. Without explicit governance mechanisms, recursive AI development
systems are prone to compounding misalignment, emergent complexity, and silent failure
modes that are difficult to detect and costly to remediate.

This paper introduces \textit{Reflector}, a framework for designing, implementing, and
governing \textit{reflective development systems}—AI-augmented engineering workflows
that incorporate structured self-evaluation, bounded autonomy, and human alignment
checkpoints as first-class architectural concerns.

\subsection{Contributions}

This paper makes the following contributions:

\begin{enumerate}
  \item A formal characterization of recursive development systems and their associated
    failure modes.
  \item A taxonomy of drift types in AI-augmented engineering workflows.
  \item The Reflector framework: four core mechanisms for governance of recursive
    development systems.
  \item An architectural reference model for implementing Reflector-compliant systems.
  \item Design guidelines for integrating Reflector into existing development pipelines.
\end{enumerate}

\subsection{Paper Organization}

The remainder of this paper is organized as follows.
Section~\ref{sec:recursive-development} characterizes recursive development systems.
Section~\ref{sec:human-in-the-loop} describes the human-in-the-loop governance model.
Section~\ref{sec:recursive-auditing} presents the recursive auditing system.
Section~\ref{sec:milestone-synchronization} defines milestone synchronization protocols.
Section~\ref{sec:governance-model} describes the governance contract model.
Section~\ref{sec:proposed-architecture} presents the Reflector architecture.
Section~\ref{sec:drift-failure-modes} characterizes drift and failure modes.
Section~\ref{sec:future-work} outlines directions for future work.
Section~\ref{sec:conclusion} concludes.
